\subsection{P7 公平 A/B 的实际结果：面积退化而非收益}
\begin{table}[H]\centering\small
\begin{tabular}{lrrr}\toprule
指标&修复后基线 P7&列共享/借位 P7&变化\\\midrule
Slice LUT&171,246&199,039&+27,793（+16.23\%）\\
FF&66,817&66,778&$-39$\\
DSP48&20&20&0\\
BRAM / latch&0 / 0&0 / 0&0\\
Setup WNS（ns）&$-17.426$&$-17.112$&+0.314\\
Hold WHS（ns）&$-0.264$&$-0.264$&0\\
Pulse width slack（ns）&+0.025&+0.025&0\\
最坏数据路径（ns）&18.941&18.627&$-0.314$\\
最坏路径逻辑级数&135&148&+13\\
该路径 CARRY4 数&117&125&+8\\
Vectorless 估计（W）&1.58&1.58&报告精度内相同\\\bottomrule
\end{tabular}
\caption{相同器件、Vivado 2018.3、P7 与申请周期1.5625\,ns的完整综合结果。均为综合后估计，尚未布局布线；功耗未使用实测开关活动。}
\end{table}

这组结果否定了“源码加法表达式变少就必然节省 FPGA 面积”的假设。
列共享与借位改写的组合使 LUT 增加16.23\%，setup只改善0.314\,ns，仍大幅违反目标。
最坏路径仍位于除法器：基线从\code{dv_reg[0]}到\code{q_reg[60]}，
改写版从\code{rem_reg[0]}到同一目标寄存器。
改写后125个CARRY4串联，说明每拍展开7步的串行借位依赖仍是主要频率障碍；
FF仅减少39个，不足以抵消组合资源增长。
因此，这一版本只能称为功能等价的算术结构候选，不能宣传已经获得低面积PPA优势。

两份\code{check_timing}的16类问题计数均为0，未约束路径表为空；
这说明记录约束下的检查覆盖，并不消除实际负setup/hold裕量。
OOC时钟尚未通过真实内部BUFG布线，且I/O延迟为零，
所以不能从18.627\,ns的综合估计直接宣布53.7\,MHz为可用频率。
后续P5以25\,ns进行FPGA筛选，另外改变了展开数、权重最高位表达式和时钟约束；
必须单列其结果。行和缓存若进入候选，也须与该P5在相同25\,ns条件下比较。

基线DCP为69,108,593字节，摘要\code{9c3bc09dd193d0bf}\ldots；
列共享/借位DCP为71,421,073字节，摘要\code{4712577f4291fed3}\ldots。
两份均已逐字节核对远端与本地SHA-256；原始报告、源归档和完整摘要见
\code{docs/evidence/20260930/vivado_complete_p7/}。
第一份运行曾在报告下载阶段超时，原失败manifest保留，随后从同一次已完成运行恢复并验证；
不能把文件传输故障混写为综合失败，也不能只凭恢复状态略去原始过程。
